% GENERATED FILE. Edit canonical lessons, then run npm run generate:books.
% canonical-source-hash: fnv1a64:56a7074a1cc366d6
% canonical-lessons: AR-C196-zaraa, AR-C196-qatafa, AR-C196-amtara, AR-C196-hannaa, AR-C196-wathiqa2

\chapter{To plant, To pick, To rain, To congratulate, To trust}
\label{ch:ar-to-plant-to-pick-to-rain-to-congratulate-to-trust}
\hlchaptermodality{196}

\begin{chapteropening}
\textbf{By the end of this chapter:} \emph{I can say to plant, to pick, to rain, to congratulate and to trust.}

Complete the last lesson of chapter 196: I can say to plant, to pick, to rain, to congratulate and to trust.
\end{chapteropening}

\section[\texorpdfstring{\ar{زرع} (zaraʿa)}{زرع (zaraʿa)}]{\ar{زرع} (zaraʿa) — to plant}

\label{lesson:AR-C196-zaraa}



\begin{quote}
Before the new one: say the Arabic for to get lost, then the Arabic for to show the way.
\end{quote}

\subsection*{You'll want to know: \ar{زرع}}
\textbf{\ar{زرع}} — \emph{zaraʿa} — \textquotedblleft{}to plant\textquotedblright{}.

\begin{cousinweb}
\textbf{\ar{زرع}} (\emph{zaraʿa}) is built on the native Arabic root \emph{z-r-ʿ}, \textquotedblleft{}to sow\textquotedblright{} (Lane). \emph{Zaraʿa} is the plain verb, \textquotedblleft{}he sowed\textquotedblright{}.
\end{cousinweb}

\begin{grammarlens}[title={What you've built}]
One more word for this chapter.
\end{grammarlens}

\subsection*{Your turn}
\begin{itemize}
\raggedright
  \item \emph{Say it:} \emph{zaraʿa}
  \item \emph{Say it:} \emph{zaraʿa}, once more
  \item \emph{Say it:} say \emph{zaraʿa}
  \item \emph{Recall:} say the Arabic for to get lost, then the Arabic for to show the way, then say \emph{zaraʿa} again
\end{itemize}

\begin{tcolorbox}[breakable,colback=teal!4,colframe=teal!35!black,title={Before you move on}]
What does \ar{زرع} mean? (\textquotedblleft{}To plant\textquotedblright{}.) Say it once more.
\end{tcolorbox}
\section[\texorpdfstring{\ar{قطف} (qaṭafa)}{قطف (qaṭafa)}]{\ar{قطف} (qaṭafa) — to pick (fruit, flowers)}

\label{lesson:AR-C196-qatafa}



\begin{quote}
Before the new one: say the Arabic for to show the way, then the Arabic for to plant.
\end{quote}

\subsection*{You'll want to know: \ar{قطف}}
\textbf{\ar{قطف}} — \emph{qaṭafa} — \textquotedblleft{}to pick (fruit, flowers)\textquotedblright{}.

\begin{cousinweb}
\textbf{\ar{قطف}} (\emph{qaṭafa}) is built on the native Arabic root \emph{q-ṭ-f}, \textquotedblleft{}to pluck, to gather\textquotedblright{} (Lane). \emph{Qaṭafa} is the plain verb, \textquotedblleft{}he plucked\textquotedblright{}.
\end{cousinweb}

\begin{grammarlens}[title={What you've built}]
One more word for this chapter.
\end{grammarlens}

\subsection*{Your turn}
\begin{itemize}
\raggedright
  \item \emph{Say it:} \emph{qaṭafa}
  \item \emph{Say it:} \emph{qaṭafa}, once more
  \item \emph{Say it:} say \emph{qaṭafa}
  \item \emph{Recall:} say the Arabic for to show the way, then the Arabic for to plant, then say \emph{qaṭafa} again
\end{itemize}

\begin{tcolorbox}[breakable,colback=teal!4,colframe=teal!35!black,title={Before you move on}]
What does \ar{قطف} mean? (\textquotedblleft{}To pick (fruit, flowers)\textquotedblright{}.) Say it once more.
\end{tcolorbox}
\section[\texorpdfstring{\ar{أمطر} (amṭara)}{أمطر (amṭara)}]{\ar{أمطر} (amṭara) — to rain}

\label{lesson:AR-C196-amtara}



\begin{quote}
Before the new one: say the Arabic for to plant, then the Arabic for to pick.
\end{quote}

\subsection*{You'll want to know: \ar{أمطر}}
\textbf{\ar{أمطر}} — \emph{amṭara} — \textquotedblleft{}to rain\textquotedblright{}.

\begin{cousinweb}
\textbf{\ar{أمطر}} (\emph{amṭara}) is built on the native Arabic root \emph{m-ṭ-r}, \textquotedblleft{}rain; to rain\textquotedblright{} (Lane). \emph{Amṭara} is Form IV, \textquotedblleft{}to rain\textquotedblright{}.
\end{cousinweb}

\begin{grammarlens}[title={What you've built}]
One more word for this chapter.
\end{grammarlens}

\subsection*{Your turn}
\begin{itemize}
\raggedright
  \item \emph{Say it:} \emph{amṭara}
  \item \emph{Say it:} \emph{amṭara}, once more
  \item \emph{Say it:} say \emph{amṭara}
  \item \emph{Recall:} say the Arabic for to plant, then the Arabic for to pick, then say \emph{amṭara} again
\end{itemize}

\begin{tcolorbox}[breakable,colback=teal!4,colframe=teal!35!black,title={Before you move on}]
What does \ar{أمطر} mean? (\textquotedblleft{}To rain\textquotedblright{}.) Say it once more.
\end{tcolorbox}
\section[\texorpdfstring{\ar{هنأ} (hannaʾa)}{هنأ (hannaʾa)}]{\ar{هنأ} (hannaʾa) — to congratulate}

\label{lesson:AR-C196-hannaa}



\begin{quote}
Before the new one: say the Arabic for to pick, then the Arabic for to rain.
\end{quote}

\subsection*{You'll want to know: \ar{هنأ}}
\textbf{\ar{هنأ}} — \emph{hannaʾa} — \textquotedblleft{}to congratulate\textquotedblright{}.

\begin{cousinweb}
\textbf{\ar{هنأ}} (\emph{hannaʾa}) is built on the native Arabic root \emph{h-n-ʾ}, \textquotedblleft{}to be wholesome, to be pleasant\textquotedblright{} (Lane). \emph{Hannaʾa} is Form II, \textquotedblleft{}to wish one joy\textquotedblright{}.
\end{cousinweb}

\begin{grammarlens}[title={What you've built}]
One more word for this chapter.
\end{grammarlens}

\subsection*{Your turn}
\begin{itemize}
\raggedright
  \item \emph{Say it:} \emph{hannaʾa}
  \item \emph{Say it:} \emph{hannaʾa}, once more
  \item \emph{Say it:} say \emph{hannaʾa}
  \item \emph{Recall:} say the Arabic for to pick, then the Arabic for to rain, then say \emph{hannaʾa} again
\end{itemize}

\begin{tcolorbox}[breakable,colback=teal!4,colframe=teal!35!black,title={Before you move on}]
What does \ar{هنأ} mean? (\textquotedblleft{}To congratulate\textquotedblright{}.) Say it once more.
\end{tcolorbox}
\section[\texorpdfstring{\ar{وثق} (wathiqa)}{وثق (wathiqa)}]{\ar{وثق} (wathiqa) — to trust}

\label{lesson:AR-C196-wathiqa2}



\begin{quote}
Before the new one: say the Arabic for to rain, then the Arabic for to congratulate.
\end{quote}

\subsection*{You'll want to know: \ar{وثق}}
\textbf{\ar{وثق}} — \emph{wathiqa} — \textquotedblleft{}to trust\textquotedblright{}.

\begin{cousinweb}
\textbf{\ar{وثق}} (\emph{wathiqa}) is built on the native Arabic root \emph{w-th-q}, \textquotedblleft{}to trust; to be firm\textquotedblright{} (Lane). \emph{Wathiqa} is the plain verb, \textquotedblleft{}he trusted\textquotedblright{}.
\end{cousinweb}

\begin{grammarlens}[title={What you've built}]
One more word for this chapter.
\end{grammarlens}

\subsection*{Your turn}
\begin{itemize}
\raggedright
  \item \emph{Say it:} \emph{wathiqa}
  \item \emph{Say it:} \emph{wathiqa}, once more
  \item \emph{Say it:} say \emph{wathiqa}
  \item \emph{Recall:} say the Arabic for to rain, then the Arabic for to congratulate, then say \emph{wathiqa} again
\end{itemize}

\begin{tcolorbox}[breakable,colback=teal!4,colframe=teal!35!black,title={Before you move on}]
What does \ar{وثق} mean? (\textquotedblleft{}To trust\textquotedblright{}.) Say it once more.
\end{tcolorbox}
